\chapter{DME均质充量压燃着火的数值模拟方法}

\section{二级标题}

正文内容

\subsection{三级标题}

正文内容

% 公式另起一行，按章编号；较长公式在等号或运算符处转行，上下式在等号处对齐
\begin{equation}
  m = \sum_{k=1}^{K} m_k
  \label{eq:sum}
\end{equation}
\begin{equation}
  \begin{aligned}
    f(x,y) &= f(0,0)+\frac{1}{1!}\Bigl(x\frac{\partial}{\partial x}+y\frac{\partial}{\partial y}\Bigr)f(0,0)\\
           &\quad +\frac{1}{2!}\Bigl(x\frac{\partial}{\partial x}+y\frac{\partial}{\partial y}\Bigr)^{2}f(0,0)+\cdots\\
           &\quad +\frac{1}{n!}\Bigl(x\frac{\partial}{\partial x}+y\frac{\partial}{\partial y}\Bigr)^{n}f(0,0)+\cdots
  \end{aligned}
  \label{eq:taylor}
\end{equation}

% 表格：三线表，表名在上方。一页放得下的表用 table + tabular（浮动体）
\begin{table}[htbp]
  \centering
  \caption{选取组分的热力学性质}\label{tab:thermo}
  \begin{tabular}{cccc}
    \toprule
    组分 & H\textsubscript{f}(kcal/mol) & S\textsubscript{f}(kcal/mol) & C\textsubscript{p}(kcal/mol) \\
    \midrule
    A1 & 100 & 100 & 100 \\
    A2 &     &     &     \\
    A3 &     &     &     \\
    A4 & 100 & 100 & 100 \\
    A5 &     &     &     \\
    A6 &     &     &     \\
    A7 &     &     &     \\
    A8 &     &     &     \\
    \bottomrule
  \end{tabular}
\end{table}

% 一页放不下、需要跨页的长表才用 longtable：续页自动重复表头，并在右上方标「续表」。
% longtable 不是浮动体，不要让它与 [t]、[htbp] 的图表挤在同一页（必要时在它前面加 \clearpage）。
% \begin{longtable}{cccc}
%   \caption{表名}\label{tab:long}\\
%   \toprule 组分 & A & B & C \\ \midrule
%   \endfirsthead
%   \sdjucontinued{4}
%   \toprule 组分 & A & B & C \\ \midrule
%   \endhead
%   \bottomrule
%   \endlastfoot
%   A1 & 100 & 100 & 100 \\
%   ……
% \end{longtable}

% 图：图名在正下方
\begin{figure}[htbp]
  \centering
  \includegraphics[width=0.55\textwidth]{fig2-1}
  \caption{气缸压力随曲轴转角变化的曲线}
  \label{fig:pressure}
\end{figure}
